% Part: many-valued-logic
% Chapter: three-valued-logics
% Section: introduction

\documentclass[../../../include/open-logic-section]{subfiles}

\begin{document}

\olfileid{mvl}{thr}{int}

\olsection{પરિચય}

$\True$ અને $\False$માં માત્ર એક વધુ મૂલ્ય~$\Undef$ ઉમેરીએ,
તો ત્રિમૂલ્ય તર્કશાસ્ત્ર મળે છે. એક જ વધારાનું સત્યમૂલ્ય
હોવા છતાં $\lnot$, $\land$, $\lor$ અને $\lif$ માટે
સત્યવિધેયો વ્યાખ્યાયિત કરવાની ઘણી રીતો છે. પછી કોઈ
તર્કશાસ્ત્ર આ સત્યવિધેયોના કોઈ પણ સંયોજનનો ઉપયોગ કરી શકે છે;
ઉપરાંત માત્ર $\True$ને નિર્દિષ્ટ ગણવું કે $\True$ અને~$\Undef$
બંનેને નિર્દિષ્ટ ગણવાં, એ પસંદગી પણ છે.

અહીં આપણે સૌથી જાણીતા ત્રિમૂલ્ય તર્કશાસ્ત્રોમાંથી કેટલાક,
તેમની પાછળની પ્રેરણા અને તેમના કેટલાક ગુણધર્મો રજૂ કરીએ છીએ.

\end{document}
